% =====================================================================
% Chapter 3, Sections 3.5.6 to 3.6, rewritten to the section formula
% (2026-10-04). Replaces everything from \subsection{Session Control}
% up to, and not including, \printbibliography.
%
% THE SECTION FORMULA
%
% \section (for example 3.6 Summary)
%   - One opening paragraph: what the section covers and how it connects
%     to the requirements or to the previous section.
%   - No mechanism detail at this level. A summary may use one mapping
%     table in place of long prose.
%
% \subsection (one component; 0.5 to 1.5 pages)
%   1. Role: one paragraph, in plain language, on what the component
%      does for the operator and which requirement it serves.
%   2. Mechanism: one to three paragraphs on how it works. Name the main
%      function or file once, where a presenter would look for it.
%   3. Decision: one paragraph on the consequential design choice and
%      its reason.
%   4. Limit or divergence: only when it changes what the reader may
%      claim. A divergence from Chapter 1 is always stated.
%   5. At most one figure or table, when a relationship or mapping is
%      clearer drawn than written.
%   When the subsection has subsubsections, it keeps only the role and a
%   short overview (at most two paragraphs, or one comparison table),
%   and the mechanisms move below it.
%
% \subsubsection (one mechanism; 0.5 to 1 page)
%   1. The first sentence states what the mechanism does.
%   2. One or two paragraphs on how it works.
%   3. One paragraph on why, or on the edge case that matters.
%   4. An optional table or figure. No code figure unless it serves as a
%      presentation anchor for the defense.
%
% Throughout: numbers that come from measurement belong in Chapter 4;
% QA history and finding numbers do not appear; identifiers appear when
% a reader needs them to find the code, and each is explained once.
%
% New labels: tab:gov-session-controls, sec:gov-kill-switch,
% fig:gov-killswitch, sec:gov-escalation-routing,
% tab:gov-escalation-settings, sec:gov-persistent-approvals,
% fig:gov-alwaysallow, sec:gov-task-slot, fig:gov-taskslot,
% sec:gov-secondary-runtime, fig:gov-codexgates, fig:gov-tenant,
% tab:gov-agent-lifecycle, tab:gov-dashboard-sections,
% tab:gov-security-limits, fig:gov-toctou, tab:gov-requirement-map.
% =====================================================================

\FloatBarrier
\subsection{Session Control}
\label{sec:gov-killswitch}

Session control lets an operator stop work that an agent has already
started. It serves Requirement~7, which asks that an agent session can be
suspended or terminated within one second. The dashboard offers two
controls with different reach, compared in
Table~\ref{tab:gov-session-controls}. \textit{Cancel} ends one prompt sent
from the dashboard. The emergency kill switch stops an entire agent and
keeps it stopped.

\begin{table}[htbp]
\centering
\caption{The two session controls.}
\label{tab:gov-session-controls}
\small
\begin{tabularx}{\textwidth}{@{}>{\raggedright\arraybackslash}p{0.2\textwidth}>{\raggedright\arraybackslash}X>{\raggedright\arraybackslash}X@{}}
\toprule
 & \textbf{Cancel} & \textbf{Emergency kill switch} \\
\midrule
What it stops & One dashboard prompt & Every running task of one agent,
and any further action by the agent or its child sessions \\
Who may use it & The account that sent the prompt, or an Administrator
or Root in the organization & A User for an assigned agent, or an
Administrator or Root for any agent in the organization \\
Later work & Unaffected. The agent can receive the next prompt & Blocked
until an operator releases the lockdown \\
Ledger & The cancellation and the account that requested it & The
lockdown, both stop times, and whether the stop was confirmed \\
\bottomrule
\end{tabularx}
\end{table}

Cancel is the ordinary control, used for a mistaken request or one that is
taking too long. The server checks it against its own record of the run,
which stores the account that started the prompt, so a caller cannot
widen its reach by naming a different agent. Canceling signals the run's
\texttt{AbortController} and marks its ending as \texttt{cancelled}. The
task then shows as \textit{Stopping} until the run has actually ended, as
Section~\ref{sec:gov-task-slot} explains. The same mechanism ends a prompt
that reaches its five-minute limit, with the ending \texttt{timeout}, so
the conversation shows whether a person or the time limit stopped it.

The kill switch is the emergency control, found in the \textit{Emergency
kill switch} section of the dashboard. It acts in two steps, lockdown and
then termination, described in the next two subsections. Viewers cannot
use either control, because both change what an agent does.

\subsubsection{Agent Lockdown}
\label{sec:gov-agent-lockdown}

Lockdown blocks every further governed action of an agent until an
operator releases it. The kill switch adds the agent's identifier to the
\texttt{lockedAgents} list in the organization's \texttt{policy.json}.
Because the list is part of the stored policy, the lockdown stays in force
after the current request ends and after the Gateway restarts.

The policy engine checks this list before it evaluates any rule, directly
after confirming that the agent is registered
(Section~\ref{sec:gov-evaluation-order}). An allow rule cannot override a
lockdown, and the check does not depend on recognizing the tool or
extracting its resource. Lockdown also blocks in the \texttt{monitor}
posture. Monitor posture relaxes ordinary rule decisions, and a lockdown is
an operator's instruction to stop. When the installation posture is
\texttt{off}, the tool-call gate does not run. The dashboard's prompt route
still checks the list in that posture and rejects a prompt to a locked
agent.

Lockdown also covers work that the agent has already handed on. A child
session can run under a different agent identifier, so the engine follows
the host's \texttt{spawnedBy} lineage for up to sixteen levels and blocks a
call whose ancestor is locked. If the lineage cannot be read while an agent
is locked, the call is blocked, because the engine cannot establish that it
is unrelated to the locked agent. The ledger records a direct lockdown, a
locked ancestor, and unreadable lineage under different rule identifiers,
so an investigation can tell them apart.

An action that was waiting for approval when the lockdown began can no
longer be allowed. Before the approval route accepts \textit{Allow once} or
\textit{Always allow}, it reads the agent's current policy and returns a
conflict response if the agent is locked. \textit{Deny} is still accepted,
so the operator can close the request.

Releasing the lockdown removes the identifier from the list and records
the operator in the ledger. Work stopped during the incident does not
restart, and held approvals are not granted. An operator starts that work
again if it is still needed.

\subsubsection{Kill Switch}
\label{sec:gov-kill-switch}

The kill switch combines lockdown with termination of the agent's running
work, in that order. \texttt{lockDownAgent} in
\texttt{src/governance/kill-switch.ts} first stores the lockdown and then
calls \texttt{terminateAgentRuns}, as Figure~\ref{fig:gov-killswitch}
shows. The order closes a gap. If running work were stopped first, the
agent could begin a new action in the interval before the lockdown was
stored.

\begin{figure}[htbp]
\centering
\begin{tikzpicture}[node distance=6mm and 6mm]
  \tikzset{ksstep/.style={gbox, text width=27mm, minimum height=17mm}}
  \node[gbox, text width=21mm, minimum height=17mm] (press)
    {Operator presses\\\textit{Lock down}};
  \node[ksstep, right=of press, draw=orange!70!black, fill=orange!11] (lock)
    {\textbf{1. Lockdown}\\\scriptsize agent added to\\\texttt{lockedAgents}};
  \node[ksstep, right=of lock, draw=red!60!black, fill=red!8] (abort)
    {\textbf{2. Termination}\\\scriptsize abort signal to every running task};
  \node[ksstep, right=of abort, draw=violet!65!black, fill=violet!9] (watch)
    {\textbf{3. Confirmation}\\\scriptsize registries checked every 10\,ms
     for up to 2\,s};
  \draw[gflow] (press) -- (lock);
  \draw[gflow] (lock) -- (abort);
  \draw[gflow] (abort) -- (watch);
  % The two measured times.
  \draw[decorate, decoration={brace, amplitude=4pt}]
    ([yshift=2mm]abort.north west) -- ([yshift=2mm]abort.north east)
    node[midway, above=3pt, font=\scriptsize] {dispatch time};
  \draw[decorate, decoration={brace, amplitude=4pt}]
    ([yshift=9mm]lock.north west) -- ([yshift=9mm]watch.north east)
    node[midway, above=3pt, font=\scriptsize] {confirmed-stop time};
  % Where the result goes.
  \node[gstore, below=9mm of abort, text width=58mm] (out)
    {Dashboard result and ledger entry \texttt{governance.agent.lock}:\\
     both times, the runs aborted, and whether the stop was confirmed};
  \draw[gflow] (watch.south) |- (out.east);
\end{tikzpicture}
\caption{The kill switch. Lockdown is stored before any running task is
signaled, and the stop is reported as confirmed only when every signaled
run has left the run registries.}
\label{fig:gov-killswitch}
\end{figure}

Termination sends an abort signal to every running task of the agent. It
reaches the prompts started from the dashboard and the host's own chat
runs, which the Gateway registers with the layer through a terminator
function. Tools that started operating-system processes end those
processes when their run is aborted, where the tool supports it.

The kill switch then reports two separate times. The dispatch time is how
long it took to send the abort signals. The confirmed-stop time runs from
the start of the operation until every signaled run has left the run
registries, which are checked every 10 milliseconds for up to two seconds.
A run still present after two seconds is reported as not confirmed
stopped. When nothing was running, the stop is reported as complete at
once. Both times, the runs aborted, and the confirmation are shown on the
dashboard and written to the ledger entry \texttt{governance.agent.lock}.

This reporting departs from the way Chapter~1 stated the target.
Requirement~7 names a one-second response, and sending the abort signal
meets it with a wide margin. The design reports the confirmed stop as
well, because a sent signal does not show that a model call or a tool has
finished. Chapter~4 gives the measured times, including runs whose
confirmed stop took longer than one second.

If the ledger entry cannot be written, the lockdown and termination still
take effect. The dashboard then shows the completed stop together with a
notice that its ledger entry is missing, so a failed audit write is not
reported as a failed emergency stop.

The dashboard keeps the kill switch available while other operations on
the page are in progress. Creating an agent is the one operation that
delays it. During creation the Gateway does not respond to any request for
about six seconds, and a stop pressed in that interval takes effect when
the Gateway responds again (Section~\ref{sec:gov-security}).

\FloatBarrier
\subsection{Human-in-the-Loop Approvals}
\label{sec:gov-approvals}

When no rule covers an action, the policy engine can ask a person before
denying it. This is the human approval that Requirement~5 records. The
engine returns an approval request, OpenClaw's approval machinery holds the
tool call, and the call continues only if an operator allows it. Only an
action that matched no rule can be escalated. An action blocked by a deny
rule or a core rule is never offered to a person.

The approval card offers three answers. \textit{Allow once} lets the
action run this time. \textit{Always allow} lets it run this time and asks
an Administrator to make the permission permanent. \textit{Deny} blocks
it. The ledger records each answer and the action's final outcome. The
following subsections describe who receives a request, and what
\textit{Always allow} leads to.

\subsubsection{Escalation Routing}
\label{sec:gov-escalation-routing}

The place where an approval request appears depends on where the run
started. A prompt sent from the governance dashboard produces a request in
the dashboard's \textit{Waiting for your answer} section, shown to the
accounts that manage the agent. These are a User assigned the agent, and
any Administrator or Root in the organization. A run started from a chat
channel produces a request on OpenClaw's own approval surfaces, the Control
UI and the channel. There, anyone holding the Gateway credential or access
to the channel can answer, as Section~\ref{sec:gov-security} notes.

When several accounts can answer one request, the first answer is used. A
later answer is rejected, and the dashboard displays a notice that it was
not used. The ledger records the answer together with the account that
gave it.

Four settings control escalation, listed in
Table~\ref{tab:gov-escalation-settings}. Where the agent setting and an
account setting both apply, the stricter one is used, so either can switch
escalation off. For a dashboard prompt the account setting is that of the
account that sent it. For a run that nobody started by name, such as a chat
message, the settings of every account assigned the agent are considered.

\begin{table}[htbp]
\centering
\caption{Escalation settings and who may change them.}
\label{tab:gov-escalation-settings}
\small
\begin{tabularx}{\textwidth}{@{}>{\raggedright\arraybackslash}p{0.27\textwidth}>{\raggedright\arraybackslash}X>{\raggedright\arraybackslash}p{0.27\textwidth}@{}}
\toprule
\textbf{Setting} & \textbf{Effect} & \textbf{Changed by} \\
\midrule
Escalation for an agent (\texttt{agentAsk}) & Ask a person when no rule
matches, or deny at once & Administrator or Root \\
Escalation for an account (\texttt{userAsk}) & The same choice, for runs
on that account's behalf & Root \\
Waiting time, installation (\texttt{hitlTimeoutSeconds}) & How long a
request waits; 300 seconds by default & Administrator or Root \\
Waiting time, one agent (\texttt{agentHitlTimeout}) & Replaces the
installation value for that agent & A User for an assigned agent,
Administrator, or Root \\
\bottomrule
\end{tabularx}
\end{table}

Chapter~1 gave the installation waiting time to Root. The built system lets
Administrators set it as well, because it sits with the other
installation-wide policy settings that Administrators control, and
Administrators are the tier that answers escalations. Every waiting time is
bounded between 5 seconds and 24 hours.

If nobody answers within the waiting time, the action is denied. The
request is kept in the \textit{Awaiting your decision} section, with how
long it waited and whether it timed out or was canceled. An operator who
manages the agent can allow or deny it there later. Allowing it does not
run the original action, which has already ended. It files the same rule
request as \textit{Always allow}, so the action can succeed when it is next
attempted. An unanswered request always ends in a denial, so an
installation with nobody watching does not fall back to allowing actions.

\subsubsection{Persistent Approvals}
\label{sec:gov-persistent-approvals}

\textit{Always allow} lets the current action run once and files a rule
request. It does not write a rule. This differs from Chapter~1, where the
answer granted permanent permission directly. A chat approval can be
answered by a person with no governance account, and a permanent change to
the policy needs an approver who is signed in and named in the ledger. The
approval card states this before the button is pressed:
``\textit{Always allow} allows this action once and requests permission for
future attempts. An Administrator must approve the request before
permission becomes permanent.''

\begin{figure}[htbp]
\centering
\begin{tikzpicture}[node distance=5mm and 5mm, every node/.append style={execute at begin node={\hyphenpenalty=10000\relax}}]
  \tikzset{
    aaAgent/.style={gbox, text width=29mm, minimum height=15mm,
      draw=teal!65!black, fill=teal!9},
    aaOperator/.style={gbox, text width=29mm, minimum height=15mm,
      draw=orange!70!black, fill=orange!11},
    aaAdmin/.style={gbox, text width=29mm, minimum height=15mm,
      draw=blue!65!black, fill=blue!8}
  }
  \node[aaAgent] (miss) {\scriptsize Tool call matches no rule; policy
    returns an approval request};
  \node[aaOperator, right=of miss] (card) {\scriptsize Card shows
    \textit{Allow once}, \textit{Always allow}, \textit{Deny}};
  \node[aaOperator, right=of card] (press) {\scriptsize Operator presses
    \textit{Always allow}};
  \node[aaAgent, right=of press] (once) {\scriptsize Action runs once};
  \node[aaAgent, below=10mm of once] (file) {\scriptsize Callback files a
    pending rule request for this resource};
  \node[aaAdmin, left=of file] (review) {\scriptsize Administrator or Root
    reviews it in \textit{Rule requests}};
  \node[aaAdmin, left=of review] (rule) {\scriptsize Approved: rule added,
    described as ``Requested by \ldots''};
  \node[aaAgent, left=of rule] (later) {\scriptsize Later calls allowed by
    the rule};
  \draw[gflow] (miss) -- (card);
  \draw[gflow] (card) -- (press);
  \draw[gflow] (press) -- (once);
  \draw[gflow] (once) -- (file);
  \draw[gflow] (file) -- (review);
  \draw[gflow] (review) -- (rule);
  \draw[gflow] (rule) -- (later);
  \node[gnote, below=2mm of file, text width=34mm] {queue full: the action
    still runs once, and a warning says the request was not saved};
  \node[gnote, below=2mm of review, text width=27mm] {rejected: no rule is
    written};
\end{tikzpicture}
\caption{\textit{Always allow}. The action is allowed once at the moment of
the answer. A permanent rule exists only after an Administrator or Root
approves the request that the answer filed. Teal marks the agent's run,
orange the operator who answered, and blue the approver.}
\label{fig:gov-alwaysallow}
\end{figure}

Figure~\ref{fig:gov-alwaysallow} shows the sequence. The rule request is
filed after the answer, by the policy's \texttt{onResolution} callback in
the agent's run. The requested rule allows exactly the resource that was
asked about, for that agent only, with the same access direction. The
request's reason is generated from the escalation, for example ``Agent
`scout' asked to run `exec' against command `uname -a'.'' When an account
answered from the dashboard, the request names that account. A request
answered on a chat surface is labeled \texttt{hitl-approval}, because no
governance account authored it.

An Administrator or Root approves or rejects the request in the
\textit{Rule requests} section. Approval writes the rule with a
description that names who asked and why, for example ``Requested by user1,
answering an escalation: Agent `scout' asked to run `exec' against command
`uname -a'.'' The rule's description then explains its origin wherever
the rule is listed. A request is decided once, and approval is rejected if the agent it
names has been deleted or replaced since the request was filed.

Requests filed by escalations share one queue per organization, sized at 40
plus 20 for every account, and an identical pending request is not filed
twice. When the queue is full, the action is still allowed once, because
the operator already granted it, and a warning states that the request was
not saved. The warning appears in the dashboard and the Control UI. Chat
channels show only the original answer.

\FloatBarrier
\subsection{Runtime Management}
\label{sec:gov-runs}

Runtime management keeps running work bounded and governed wherever the
agent executes. Prompts sent from the dashboard run as tasks with time and
concurrency limits, described in Section~\ref{sec:gov-task-slot}. OpenClaw
can also run an agent inside a separate helper process, and
Section~\ref{sec:gov-secondary-runtime} describes how the gate reaches that
runtime and who may permit it.

\subsubsection{Task and Slot Model}
\label{sec:gov-task-slot}

The run registry in \texttt{src/governance/prompt-runs.ts} tracks every
dashboard prompt as a task. A task has two parts that end at
different moments. One is the row an operator sees, and the other is the
slot the task occupies under the concurrency limits. An account may have two prompts running at once,
the installation six, and each prompt may run for five minutes. A prompt
beyond either concurrency limit is rejected with a message that names the
limit.

\begin{figure}[htbp]
\centering
\begin{tikzpicture}[node distance=6mm and 13mm]
  \tikzset{tsState/.style={gbox, text width=24mm, minimum height=16mm}}
  \node[tsState, draw=teal!65!black, fill=teal!9] (run)
    {\textbf{Running}\\\scriptsize listed, holds a slot,\\Cancel available};
  \node[tsState, right=of run, draw=orange!70!black, fill=orange!11] (stop)
    {\textbf{Stopping}\\\scriptsize listed, holds a slot,\\Cancel disabled};
  \node[tsState, right=of stop, draw=blue!65!black, fill=blue!8] (save)
    {\textbf{Saving reply}\\\scriptsize listed,\\slot released};
  \node[tsState, right=of save] (done)
    {\textbf{Finished}\\\scriptsize no longer listed};
  \draw[gflow] (run) -- node[above, font=\scriptsize] {stop} (stop);
  \draw[gflow] (stop) -- node[above, font=\scriptsize] {ends} (save);
  \draw[gflow] (save) -- node[above, font=\scriptsize] {saved} (done);
  \draw[gflow] (run.south) -- ++(0,-6mm) -| node[glab, pos=0.25]
    {model returns} (save.south);
  \draw[decorate, decoration={brace, amplitude=4pt}]
    ([yshift=2mm]run.north west) -- ([yshift=2mm]stop.north east)
    node[midway, above=3pt, font=\scriptsize]
    {counts toward 2 per account and 6 per installation};
\end{tikzpicture}
\caption{The states of a dashboard task. A stop can come from Cancel, the
five-minute limit, or the kill switch. The slot is released when the run
stops executing, and the row is removed only when the reply and its ledger
entry have been saved.}
\label{fig:gov-taskslot}
\end{figure}

Figure~\ref{fig:gov-taskslot} shows the states. A task appears in its
conversation and in the \textit{Active agent sessions} section, where its
owner, an Administrator, or Root can cancel it. A request to stop moves it
to \textit{Stopping}. It keeps its slot in that state, because an abort
signal does not show that the model call or the tool has finished. When
the run ends, the task moves to \textit{Saving reply} and releases its
slot. Its row stays until the reply and the ledger entry are saved, so a
task that leaves the list has finished and has been recorded.

The two parts are released separately for a reason. A slot held while the
reply was being saved would count work that had already stopped executing,
and a quick series of prompts from one account would be rejected for
capacity that was not in use. The limits count only work that is
executing.

\subsubsection{Secondary Runtime Governance}
\label{sec:gov-secondary-runtime}

OpenClaw can run an agent in two arrangements, introduced in
Figure~\ref{fig:gov-twopaths}. In the in-process arrangement, every tool
call passes through \texttt{runBeforeToolCallHook} inside the Gateway. In
the native-harness arrangement, used by the Codex app-server backend, a
separate helper process executes the tools itself. The host then reaches
the gate through a relay hook that it writes into the helper's
configuration when the session starts. Before each tool call the helper
runs the relay command, which calls back into the host, and the host runs
the same hook and returns its result.

The layer makes this relay a requirement. \texttt{governanceRequiresNativeToolRelay}
in \texttt{src/governance/native-relay-requirement.ts} makes the host
install the relay for every tool on a governed installation. If the relay
cannot reach the host, the helper receives a block response. A governed
installation has no configuration in which the helper runs a tool without
consulting the gate.

The native harness has one weakness the gate cannot remove. A recursive
search runs inside the helper and returns its results directly to the
model, so a result naming a denied file is recorded in the ledger and still
reaches the model. On the in-process runtime the layer removes such
results. Because of this gap, two permissions control whether an agent may
use the native harness, as Figure~\ref{fig:gov-codexgates} shows.

\begin{figure}[htbp]
\centering
\begin{tikzpicture}[node distance=6mm and 8mm]
  \node[gbox, text width=22mm] (session) {Agent session on the Codex
    backend};
  \node[gbox, right=of session, text width=30mm, draw=violet!65!black,
    fill=violet!9] (backend) {Backend enabled for the installation?\\
    \scriptsize set by Root};
  \node[gbox, right=of backend, text width=30mm, draw=blue!65!black,
    fill=blue!8] (agent) {Agent permitted?\\\scriptsize set by its owner
    or Root};
  \node[gbox, right=of agent, text width=24mm, draw=green!55!black,
    fill=green!10] (relay) {Tool calls relayed to the gate};
  \node[gbox, below=8mm of backend, text width=30mm] (noBackend)
    {\scriptsize Backend unavailable; agents run in-process};
  \node[gbox, below=8mm of agent, text width=40mm, draw=red!60!black,
    fill=red!8] (noAgent) {\scriptsize Each tool call blocked and recorded
    as \texttt{agent-not-permitted-on-codex}};
  \draw[gflow] (session) -- (backend);
  \draw[gflow] (backend) -- node[glab, above=1mm] {yes} (agent);
  \draw[gflow] (agent) -- node[glab, above=1mm] {yes} (relay);
  \draw[gflow] (backend) -- node[glab, right] {no} (noBackend);
  \draw[gflow] (agent) -- node[glab, right] {no} (noAgent);
\end{tikzpicture}
\caption{The two permissions for the native Codex harness. Both are off
until someone turns them on, and an agent needs both. The in-process
runtime needs neither.}
\label{fig:gov-codexgates}
\end{figure}

Root enables or disables the Codex backend for the whole installation,
because the switch changes OpenClaw's own configuration and affects model
access for every agent. The Administrator who owns an agent, or Root,
decides whether that agent may use the backend, which is stored as
\texttt{codexAllowed} in its registry record. When a tool call arrives
through the native relay from an agent without that permission, the policy
engine blocks it and records \texttt{agent-not-permitted-on-codex}. The
block also applies in the \texttt{monitor} posture, because it concerns
where the agent may run at all. The check comes after the lockdown check,
so a locked agent's calls are recorded as a lockdown.

\FloatBarrier
\subsection{Tenancy and Agent Registry}
\label{sec:gov-tenant}

An organization is the top-level governance boundary. It groups the
installation's accounts, registered agents, policy, and audit records under
one Root account. Every account and agent record carries the
organization's identifier, \texttt{groupId}, and the organization's own
files are stored in a folder named after it, as
Figure~\ref{fig:gov-tenant} shows.

\begin{figure}[htbp]
\centering
\begin{tikzpicture}[node distance=6mm and 9mm]
  \node[gstore, align=left, font=\scriptsize] (home)
    {\textbf{Governance directory} (installation-wide)\\[3pt]
     \begin{tabular}{@{}ll@{}}
     \texttt{users.json} & accounts, each with \texttt{groupId}\\
     \texttt{agents.json} & agents, each with \texttt{groupId}\\
     \texttt{sessions.json} & dashboard sessions\\
     \texttt{ledger.key} & installation key\\
     \texttt{ledger-checkpoint.json} & ledger checkpoints
     \end{tabular}};
  \node[gstore, right=of home, align=left, font=\scriptsize] (org)
    {\textbf{\texttt{groups/}\textit{groupId}\texttt{/}} (one organization)\\[3pt]
     \begin{tabular}{@{}ll@{}}
     \texttt{policy.json} & rules, postures, lockdowns\\
     \texttt{audit-ledger.jsonl} & the audit ledger\\
     \texttt{rule-requests.json} & rule requests\\
     \texttt{pending-decisions.json} & unanswered escalations\\
     \texttt{conversations.json} & conversations and attachments
     \end{tabular}};
  \draw[gflow] (home.east) -- (org.west);
\end{tikzpicture}
\caption{Where governance state is stored. Records in the installation-wide
files carry the organization's identifier, which selects the folder holding
that organization's policy, ledger, and requests.}
\label{fig:gov-tenant}
\end{figure}

One installation holds one organization, as decided in
Section~\ref{sec:gov-tenancy-model}. Creating a second Root on a claimed
installation is rejected, and a second organization needs a second
installation. The organization identifier still does work within that one
organization. Each route that acts on an agent checks that the agent is
registered to the caller's organization, and the same check rejects an
identifier that was never registered or has been deleted. When an
organization is deleted, its audit ledger is kept in its own folder, apart
from the organization created after it. Separation between two
organizations that exist at the same time is outside the supported
deployment.

The agent registry, \texttt{agents.json}, gives the layer its own record
of each agent. A record holds the host's agent identifier, a display name,
the organization, the owning Administrator, the creation time, and the
Codex permission. Registration is required for an agent to act. The policy
engine blocks every tool call from an agent with no registry record,
because the engine cannot tell which organization's policy applies to it.
The same requirement keeps the ownership rules complete, since an agent
nobody registered can neither act nor be assigned to an account.

The registry also separates agents' workspaces. OpenClaw places each later
agent's workspace inside the default agent's workspace folder.
\texttt{nestedAgentWorkspaceRoots} in
\texttt{src/governance/agent-workspace-roots.ts} lists those nested
workspaces, and the policy engine treats a path inside one as outside the
current agent's workspace. The comparison uses fully resolved paths, so a
link or another spelling of the path leads to the same result. A read
there is escalated or denied, and in-process searches remove results from
it. This keeps a User from reading, through the default agent, the files
of an agent assigned to someone else.

\subsection{Agent Lifecycle}
\label{sec:gov-lifecycle}

An agent passes through registration, changes of ownership, and finally
removal or deletion. Table~\ref{tab:gov-agent-lifecycle} lists each
operation, who may perform it, and what it changes. Every operation is
recorded in the ledger, and none of them alters the ledger's existing
entries.

\begin{table}[htbp]
\centering
\caption{Agent lifecycle operations.}
\label{tab:gov-agent-lifecycle}
\small
\begin{tabularx}{\textwidth}{@{}>{\raggedright\arraybackslash}p{0.2\textwidth}>{\raggedright\arraybackslash}p{0.18\textwidth}>{\raggedright\arraybackslash}X@{}}
\toprule
\textbf{Operation} & \textbf{Who} & \textbf{Effect} \\
\midrule
Create an agent & Administrator or Root & Adds the agent to OpenClaw's
configuration and registers it, with an owner \\
Register an existing agent & Administrator or Root & Registers an agent
that OpenClaw already has \\
Rename & Owner or Root & Changes the display name. The identifier is
fixed because rules and ledger entries use it \\
Change owner & Owner or Root & Moves the agent to another Administrator or Root.
Users and Viewers of the previous owner lose the agent \\
Remove from governance & Owner or Root & Deletes the registry record and
all assignments. The agent stays in OpenClaw and can no longer act \\
Delete from OpenClaw's agent list & Owner or Root & Also deletes the
agent's configuration entry and clears its governance policy. Its files,
jobs, and sessions stay on the host \\
Delete as OpenClaw does & Owner or Root & Runs OpenClaw's own deletion,
which moves the agent's files to the trash, and clears its governance
policy and conversations \\
\bottomrule
\end{tabularx}
\end{table}

An agent is created or registered by an Administrator or Root. The
dashboard offers registration for the agents that OpenClaw has.
Identifiers that appear only in policy rules are listed as named in the
policy, with no control to register them. If the identifier already carries rules or
settings in the policy, the registration result lists them, so the
operator sees what the new agent inherits.

Ownership determines who administers an agent. The dashboard offers the
change of owner to Root, which can list every Administrator, and the
confirmation states before the change that the previous owner's Users and
Viewers will lose the agent. An Administrator that
still owns agents cannot be demoted or deleted until each of its agents has
a new owner, so an agent is never left with an owner who cannot manage it.

Removing an agent from governance deletes only the governance record. The
agent stays in OpenClaw, and because registration is required, its tool
calls are blocked from then on. Its rules stay in the policy and apply
again if it is registered later.

Deleting an agent offers two choices, and neither is preselected. Deleting
it from OpenClaw's agent list removes its configuration entry, and the
host keeps its scheduled jobs, approvals, session records, and files. A new
agent created under the same name inherits them, and its creation result
lists what it inherited. Deleting it as OpenClaw does runs OpenClaw's own
\texttt{agents.delete}, which removes those items and moves the agent's
files to the trash folder in the Gateway account's home directory. The
layer also removes the agent's dashboard conversations and unsent
attachments in that case, keeping any attachment that a ledger entry names.
The layer runs OpenClaw's operation through a registered function and does
not reimplement it, so the host's own safety checks apply. If OpenClaw's
deletion returns an error, for example because the agent is the default
agent or is busy, the error is returned before the layer changes anything, and the
agent stays registered and governed.

Deleting the organization asks the same two-way question once, for every
agent. The organization's audit ledger is kept.

\subsection{Management Interface}
\label{sec:gov-dashboard}

Operators manage the layer through two surfaces. The first is the set of
governance HTTP routes served by the Gateway. The second is the governance
dashboard in the Control UI, which calls only those routes. Every action has a single server-side
permission check, in its route. The dashboard is built to show a section or
a control only to an account that the matching route accepts, so that a
visible control is one that can succeed. Table~\ref{tab:gov-dashboard-sections}
groups the sections by purpose.

\begin{table}[htbp]
\centering
\caption{Dashboard sections, grouped by purpose.}
\label{tab:gov-dashboard-sections}
\small
\begin{tabularx}{\textwidth}{@{}>{\raggedright\arraybackslash}p{0.17\textwidth}>{\raggedright\arraybackslash}p{0.33\textwidth}>{\raggedright\arraybackslash}X@{}}
\toprule
\textbf{Purpose} & \textbf{Sections} & \textbf{What an operator does
there} \\
\midrule
Working with agents & \textit{Your agents}; the organization's agent
registry & Prompt an agent and read its conversation. Administrators and
Root create, register, edit, and delete agents \\
Oversight & \textit{Waiting for your answer}; \textit{Awaiting your
decision}; \textit{Active agent sessions}; \textit{Emergency kill switch} &
Answer escalations, cancel tasks, and lock down agents \\
Policy & \textit{Policy}; \textit{Agent permissions}; \textit{Rule
requests} & Read and write rules, grant folders, set postures and
escalation, and propose or decide rule requests \\
Audit & \textit{Audit ledger} & Filter entries and verify the chain \\
Accounts & \textit{Identity}; \textit{Accounts} (Root); \textit{Your
accounts} (Administrator) & Sign in, and manage accounts and their
assigned agents \\
Installation & \textit{Deployment and network posture}; the organization
settings (Root) & Check the deployment and delete the organization \\
\bottomrule
\end{tabularx}
\end{table}

The policy views title every rule by its description and show its
regular expression underneath. Adding a rule requires a description of up
to 500 characters, and a folder grant requires a purpose. When a new rule
is broader than it looks, for example a pattern that matches every command
or one that lets the agent run an interpreter with code of its own
choosing, a warning is shown when the rule is added and when an
Administrator previews a rule request that would create it. The warning
does not block the rule.

The deployment report is Root's checklist of the installation's security
state. Its first row states the posture in force. ``Governance is
enforcing'' fails when governance is off, warns when the installation or
any agent is in \texttt{monitor} and names them, and passes only when every
agent enforces. Further rows report whether the ledger file is protected
from rewriting in place and whether any ledger integrity alert is still
unacknowledged. Administrators and Root also see each such alert at the
top of the dashboard until Root acknowledges it with a written reason.

On an oversight surface the displayed text is part of the control,
because an operator acts on what the page shows. A rejected action shows
its reason and the next step to take. A section that could not be
refreshed is marked as possibly out of date. The \textit{Audit ledger}
section loads the latest 200 entries and shows 50 at first, and
verification always reads the whole chain. The browser stores no
governance data. The Control UI's service worker caches only build files,
and every governance response is sent with \texttt{Cache-Control:
no-store}.

\subsection{System Security}
\label{sec:gov-security}

This subsection collects the security boundaries of the design and the
limits that apply within them. Table~\ref{tab:gov-security-limits}
summarizes them, and the two that need more explanation follow the table.

\begin{table}[htbp]
\centering
\caption{Security boundaries and their limits.}
\label{tab:gov-security-limits}
\small
\begin{tabularx}{\textwidth}{@{}>{\raggedright\arraybackslash}p{0.16\textwidth}>{\raggedright\arraybackslash}X>{\raggedright\arraybackslash}X@{}}
\toprule
\textbf{Area} & \textbf{Protection} & \textbf{Limit} \\
\midrule
Reaching the dashboard & The Gateway listens on the local machine and is
reached through an SSH tunnel. A claimed installation rejects a second
Root & An unclaimed installation is claimed by whoever creates Root first,
so Root is created immediately after deployment \\
File paths & The tool opens the exact path the gate evaluated
(Figure~\ref{fig:gov-toctou}) & A target replaced after the check, or a
link created at a path that did not yet exist, is not detected \\
Search results & In-process searches remove results the policy denies &
On the native Codex harness such results are recorded and still reach the
model (Section~\ref{sec:gov-secondary-runtime}) \\
Prompt injection & Every action a persuaded agent attempts is checked and
recorded & The layer limits what an agent does. It does not stop the
agent from being persuaded \\
Interpreter rules & A warning when a rule allows running arbitrary code &
Such a rule can be used to edit the policy file, as explained below \\
Audit ledger & Keyed hash chain, checkpoint, integrity alerts, and the
dashboard witness (Section~\ref{sec:gov-ledger-verification}) & The key,
checkpoint, and ledger are on one host. A holder of the key can forge
entries and alerts. In-place append-only protection exists only on Windows
for a Gateway without administrator rights \\
Chat approvals & \textit{Always allow} only files a rule request & Anyone
with the Gateway credential or channel access can answer a chat run's
approval \\
Kill switch & Available during every other dashboard operation & While an
agent is being created, the Gateway does not respond for about six seconds
\\
Tool coverage & Every host tool is either governed or listed with a reason
(Table~\ref{tab:gov-ungoverned}) & Ungoverned tools run without a policy
check \\
\bottomrule
\end{tabularx}
\end{table}

The first explained limit concerns the time between a check and the use
of a file. The gate resolves the path that the agent named, following any
symbolic links, evaluates the resolved file, and passes that resolved path
to the tool in place of the original string. Repointing a link after the
check has no effect, because the tool does not resolve the link
again. Figure~\ref{fig:gov-toctou} shows this case and the case that is
still open. If the resolved file itself is replaced by a link after the
check, the tool opens the new link. The same applies when the path did not
exist at the time of the check and a link is created there before the tool
writes.

\begin{figure}[htbp]
\centering
\begin{tikzpicture}[node distance=4mm and 5mm]
  \tikzset{
    tStep/.style={gbox, text width=33mm, minimum height=14mm, font=\scriptsize},
    tRow/.style={font=\small\bfseries, align=right, text width=22mm}
  }
  % Row 1: the closed case.
  \node[tRow] (r1) {Link\\repointed\\\normalfont\scriptsize(closed)};
  \node[tStep, right=of r1] (a1) {Gate resolves \texttt{notes} to
    \texttt{report.txt} and allows it};
  \node[tStep, right=of a1] (a2) {Link \texttt{notes} repointed to a
    secret file};
  \node[tStep, right=of a2, draw=green!55!black, fill=green!10] (a3) {Tool
    opens \texttt{report.txt}, the path the gate evaluated};
  % Row 2: the open case.
  \node[tRow, below=8mm of r1] (r2) {Target\\replaced\\\normalfont\scriptsize(open)};
  \node[tStep, right=of r2] (b1) {Gate resolves \texttt{notes} to
    \texttt{report.txt} and allows it};
  \node[tStep, right=of b1] (b2) {\texttt{report.txt} replaced by a link
    to a secret file};
  \node[tStep, right=of b2, draw=red!60!black, fill=red!8] (b3) {Tool opens
    \texttt{report.txt} and follows the new link};
  \foreach \x/\y in {a1/a2, a2/a3, b1/b2, b2/b3} { \draw[gflow] (\x) -- (\y); }
  \draw[gflow, black!45] ([yshift=-6mm]b1.south west) --
    node[below, font=\scriptsize] {time} ([yshift=-6mm]b3.south east);
\end{tikzpicture}
\caption{The time between checking a path and opening it. Passing the
evaluated path to the tool closes the first case. The second case is still
open.}
\label{fig:gov-toctou}
\end{figure}

The second explained limit concerns rules that allow an interpreter. The
command denials that protect the governance directory match the text of a
command, so they cannot recognize a path that the command assembles while
it runs. An allowance such as \texttt{\textasciicircum python3 .*\$} lets the agent pass
its own code to an interpreter, and that code can build the location of
\texttt{policy.json} from fragments that match no denial. Tools run under
the same operating-system account as the Gateway, so the code can then set
its own agent to \texttt{monitor} or the installation to \texttt{off}. The
ledger records the interpreter command and no posture change, because the
edit does not pass through the governance functions that record
administrative actions. \texttt{describeRuleRisks} in
\texttt{src/governance/rule-validation.ts} returns the
\texttt{runs-arbitrary-code} warning for such a rule. The warning does not
block the rule, because some agents need to run code for their work.
Closing the exposure fully would require running tools under a separate
operating-system account that cannot reach the governance directory. That
arrangement has not been evaluated with the layer and is left as future
work in Chapter~5.

\FloatBarrier
\section{Summary}
\label{sec:gov-designsummary}

This chapter described the design of the governance layer, from the
requirements in Section~\ref{sec:gov-requirements} through the
alternatives considered to the components that were built. Every tool call
passes one gate, which applies the organization's rules in a fixed order
and denies by default. Every decision and administrative action is written
to a keyed, hash-chained ledger. Four account tiers determine who may change
the policy and who may act on each agent, and operators intervene through
cancellation, lockdown, approvals, and the dashboard.
Table~\ref{tab:gov-requirement-map} maps each requirement to the part of
the design that addresses it.

\begin{table}[htbp]
\centering
\caption{Where the design addresses each requirement.}
\label{tab:gov-requirement-map}
\small
\begin{tabularx}{\textwidth}{@{}p{0.06\textwidth}>{\raggedright\arraybackslash}X>{\raggedright\arraybackslash}p{0.3\textwidth}@{}}
\toprule
\textbf{No.} & \textbf{Requirement, in brief} & \textbf{Addressed in} \\
\midrule
1 & Node.js and TypeScript with static type checking &
Section~\ref{sec:gov-arch} \\
2 & Secure dashboard with role-based access control &
Sections~\ref{sec:gov-rbac} and~\ref{sec:gov-dashboard} \\
3 & Default-deny, policy-based restriction of agent access &
Section~\ref{sec:gov-gate} \\
4 & Fine-grained and time-limited privileges & Section~\ref{sec:gov-gate}
\\
5 & Recording of every action, decision, and approval &
Sections~\ref{sec:gov-ledger} and~\ref{sec:gov-approvals} \\
6 & Tamper-evident audit logging &
Section~\ref{sec:gov-ledger-verification} \\
7 & Suspending or terminating an agent within one second &
Section~\ref{sec:gov-killswitch} \\
8 & No plaintext secrets in log files &
Section~\ref{sec:gov-ledger-sanitization} \\
9 & Deployment on Linux with open-source components &
Section~\ref{sec:gov-constraints} \\
\bottomrule
\end{tabularx}
\end{table}

The design states its limits alongside its protections, as
Section~\ref{sec:gov-security} sets out. Chapter~4 tests the completed
design against each requirement and reports the measurements, including
the kill-switch timings and the ledger's tamper detection.
